\section{Taller resuelto: Dinámica y trabajo giroscópico}

\noindent\textbf{Nivel:} Intermedio.

\apartado{1. Enunciado}{Para el tensor Ixx=9800, Iyy=14800, Izz=19200 e Ixz=760 kg m², omega=(0.13,0.05,-0.09) rad/s y momentos (600,-380,420) N m, calcule aceleraciones y tasa de energía de rotación.}

\apartado{2. Datos}{El tensor usa -Ixz en entradas xz y zx y es constante en cuerpo.}

\apartado{3. Incógnitas}{Las magnitudes y comprobaciones solicitadas en el caso.}

\apartado{4. Sistema de referencia}{Cuerpo y momentos respecto al CG; producto integral Ixz positivo según dato.}

\apartado{5. Hipótesis físicas}{Sólido rígido y masa constante en el intervalo dinámico.}

\apartado{6. Modelo matemático}{Dinámica y trabajo giroscópico}

\apartado{7. Ecuaciones simbólicas}{Se parte de los balances y definiciones siguientes.}

\begin{equation}I\dot{\vect\omega}=\vect M-\vect\omega\times I\vect\omega\end{equation}

\begin{equation}\dot E_r=\vect\omega^TI\dot{\vect\omega}=\vect\omega^T\vect M\end{equation}

\apartado{8. Desarrollo algebraico}{Las operaciones se realizan en este orden, conservando las condiciones del modelo.}

\begin{enumerate}

\item Se calcula el momento angular antes del término giroscópico.

\item Se resuelve el sistema con el tensor completo.

\item Se evalúa la tasa de energía por dos vías: omega por M y omega por I odot.

\item El producto giroscópico no aporta trabajo porque es perpendicular a omega.

\end{enumerate}

\apartado{9. Sustitución numérica}{Se conservan los resultados intermedios sin redondear y se muestran todas las operaciones escalares disponibles.}

\textbf{Momento angular, producto vectorial y sistema de ecuaciones.} Primero se calculan las tres filas de $\boldsymbol H=I\boldsymbol\omega$. Las entradas $xz$ y $zx$ son negativas con el convenio del enunciado.\[\fitmath{\boldsymbol H=\left[\begin{matrix}1342.4\\740.0\\-1826.8\end{matrix}\right]\,\mathrm{kg\,m^2/s}}\]\[\fitmath{H_{1}=(9800)(0.13)+(-760)(-0.09)\approx1342.4\,\mathrm{kg\,m^2/s}}\]\[\fitmath{H_{2}=(14800)(0.05)\approx740\,\mathrm{kg\,m^2/s}}\]\[\fitmath{H_{3}=(-760)(0.13)+(19200)(-0.09)\approx-1826.8\,\mathrm{kg\,m^2/s}}\]\[\fitmath{c_{1}=(0.05)(-1826.8)-(-0.09)(740)\approx-24.74\,\mathrm{N\,m}}\]\[\fitmath{c_{2}=(-0.09)(1342.4)-(0.13)(-1826.8)\approx116.668\,\mathrm{N\,m}}\]\[\fitmath{c_{3}=(0.13)(740)-(0.05)(1342.4)\approx29.08\,\mathrm{N\,m}}\]El producto giroscópico se desarrolla como $(qH_z-rH_y,\ rH_x-pH_z,\ pH_y-qH_x)$; su valor es:\[\fitmath{\boldsymbol\omega\times\boldsymbol H=\left[\begin{matrix}-24.74\\116.668\\29.08\end{matrix}\right]\,\mathrm{N\,m}}\]Se resta de los momentos aplicados, obteniendo el segundo miembro $\boldsymbol b=\boldsymbol M-\boldsymbol\omega\times\boldsymbol H$.\[\fitmath{\boldsymbol b=\left[\begin{matrix}624.74\\-496.668\\390.92\end{matrix}\right]\,\mathrm{N\,m}}\]La ecuación de $y$ queda desacoplada: $\dot q=b_y/I_{yy}$. Para $x,z$, eliminar la otra aceleración en el sistema $I_{xx}\dot p-I_{xz}\dot r=b_x$, $-I_{xz}\dot p+I_{zz}\dot r=b_z$ da:\[\fitmath{\dot p=\frac{I_{zz}b_x+I_{xz}b_z}{I_{xx}I_{zz}-I_{xz}^2},\qquad\dot r=\frac{I_{xz}b_x+I_{xx}b_z}{I_{xx}I_{zz}-I_{xz}^2}}\]\[\fitmath{(\dot p,\dot q,\dot r)=\left[\begin{matrix}0.0655291072083522 & -0.0335586486486487 & 0.0229542771603306\end{matrix}\right]\,\mathrm{s^{-2}}}\]Para comprobar, multiplicar el vector de aceleraciones por $I$ y sumar el término giroscópico recupera $\boldsymbol M$. El trabajo giroscópico se anula por perpendicularidad; la tasa de energía es $\boldsymbol\omega\cdot\boldsymbol M=21.2\,\mathrm W$.





\apartado{10. Resultados con unidades}{Los resultados pertenecen al modelo y al intervalo de datos declarados.}

\begin{center}\begin{tabular}{lr}\toprule Magnitud&Resultado\\\midrule

pdot & $0.06552911\,\mathrm{s^{-2}}$\\

qdot & $-0.03355865\,\mathrm{s^{-2}}$\\

rdot & $0.02295428\,\mathrm{s^{-2}}$\\

Tasa de energía & $21.2\,\mathrm W$\\

\bottomrule\end{tabular}\end{center}

\apartado{11. Interpretación}{La energía puede aumentar o disminuir según omega por M, aun cuando algunas tasas angulares tengan signo distinto.}

\apartado{12. Comprobación independiente y alcance}{Residuo vectorial y igualdad de las dos tasas energéticas, ejecutados con álgebra matricial.}

\apartado{13. Limitaciones}{No se interpreta un literal de aceleración como prueba del balance; las operaciones y residuos están conservados en el programa.}

\par\smallskip\noindent{\footnotesize\textit{Datos y solución originales. Las sustituciones registradas se reproducen con SymPy y Studium; eso no establece validación experimental ni revisión experta.}}
